\section{Experimental Setup}
\label{sec:setup}

\subsection{Environment}
\label{subsec:env}
\begin{itemize}[leftmargin=*,itemsep=0pt,topsep=1pt]
  \item \textbf{Cloud}: AWS EC2 \texttt{m5.large} ($2$\,vCPU,
        $8$\,GiB RAM, \texttt{us-west-2}), Ubuntu~22.04~LTS.
  \item \textbf{Cluster}: K3s v1.30 (single-node, embedded SQLite
        data store), Helm~3, \texttt{metrics-server} v0.7, and
        KEDA~v2.13 installed via its official Helm chart.
  \item \textbf{Broker}: a repository-owned RabbitMQ
        \texttt{StatefulSet} using image \texttt{rabbitmq:3.12-management},
        with \texttt{limits.cpu}\,=\,$1\,000$\,m,
        \texttt{limits.memory}\,=\,$1$\,GiB, and
        \texttt{tcpSocket:5672} readiness/liveness probes to avoid
        false failures on the constrained node.
\end{itemize}

\subsection{Workload}
\label{subsec:workload}
The pattern defined in
\texttt{load-test/patterns/burst.yaml} consists of three phases:
a $10$\,s warm-up at $5$\,messages/s (so that the collector captures
a steady-state baseline at \texttt{replicas}\,=\,1), a $10$\,s burst
\emph{targeting} $1\,000$\,messages/s, and $180$\,s of silence. The
harness script \texttt{run-experiment.sh} also waits $10$\,s before
launching the producer so the collector can record a baseline, and
then appends a further $180$\,s post-burst observation window. Each
trial therefore lasts $\approx$\,$390$\,s wall-clock, with roughly
$355$--$360$\,s of observation after the burst ends.

\paragraph{Producer ceiling.}~The
\texttt{pika}~\cite{pika2024} \texttt{BlockingConnection.basic\_publish}
call is a synchronous RPC that saturates at $\approx$\,$740$\,msg/s on
the \texttt{m5.large} node; consequently, the realised burst injects
$7\,100$--$7\,500$ messages rather than the nominal $10\,000$. The
ceiling affects both groups identically---the same producer publishes
to both queues with the same pattern---so it does not bias the
comparison.

\subsection{Metrics and Collection}
\label{subsec:metrics}
The collector \texttt{scripts/collect-metrics.sh} polls
\texttt{kubectl} for pod readiness and the RabbitMQ Management API
for queue statistics, emitting an 8-column CSV (see
artifact appendix). We derive three metric axes:

\begin{itemize}[leftmargin=*,itemsep=0pt,topsep=1pt]
  \item \textbf{Responsiveness.} Reaction latency (burst start
        $\rightarrow$ first new pod Ready), scale-up time (burst start
        $\rightarrow$ peak pod count), and drain time (burst end
        $\rightarrow$ \texttt{queue\_depth}\,=\,0).
  \item \textbf{User-visible effect.} Average delivery rate during the
        drain window and total messages delivered per trial.
  \item \textbf{Resource utilisation.} Scale-down-start time (burst end
        $\rightarrow$ first pod removal), pods still Ready at the end
        of the trial ($\approx$\,$355$\,s post-burst), and
        pod-seconds overshoot, defined as
        $\int \max\{0,\,\mathrm{pods}(t) - \mathrm{minReplicas}\}\,dt$
        over the entire trial.
\end{itemize}

\subsection{Trial Protocol}
\label{subsec:protocol}
For each of the two groups we run $n\,=\,3$ independent trials and
report the sample mean with one sample standard deviation. Between
trials the Deployment is scaled back to $1$ replica and the target
queue is purged. Each archived raw CSV therefore corresponds to the
same procedure: reset deployment, purge queue, capture a baseline,
run the burst pattern, then observe the drain and any scale-down.
All six formal trials were executed sequentially on 2026-04-20 and the
committed artefacts under \texttt{results/raw/} are the ground truth
for the statistics reported in this paper.
